\documentclass[11pt]{article}
\usepackage[margin=1in]{geometry}
\usepackage{amsmath,amssymb,amsthm}
\usepackage{graphicx}
\usepackage[colorlinks=true,linkcolor=blue,citecolor=blue,urlcolor=blue]{hyperref}

\newtheorem{theorem}{Theorem}
\newtheorem{lemma}{Lemma}
\newtheorem{remark}{Remark}

\newcommand{\R}{\mathbb R}
\newcommand{\one}{\mathbf 1}
\newcommand{\Linfty}{L_\infty[0,1]}
\newcommand{\Lone}{L_1[0,1]}

\title{A Counterexample to the Reverse Exact-Dominant Narrowness Question}
\author{Automatic Math Discovery Smoke Test}
\date{July 1, 2026}

\begin{document}
\maketitle

\section*{Source and Extracted Question}

The source is Nariman Abasov, Abd El Monem Megahed, and Marat Pliev,
\emph{Dominated operators from a lattice-normed space to a sequence vector
lattice}, arXiv:1508.03275.  In Section 5 the authors prove that, under their
Banach--Kantorovich hypotheses, order narrowness of a dominated linear
operator \(T\) implies order narrowness of its exact dominant \(|T|\).  The
following remark asks whether the converse is true:

\begin{quote}
Does the order narrowness of the operator \(|T|\) imply the order narrowness
of \(T\)?
\end{quote}

\begin{center}
\includegraphics[width=\textwidth]{figures/open_problem_crop.png}
\end{center}

\section*{Counterexample}

Let the measure space be the Lebesgue interval \(([0,1],\mathcal L,m)\).
Put
\[
  E=\Linfty,\qquad V=\Linfty,\qquad p(f)=|f|.
\]
Thus \(V\) is the usual vector lattice \(L_\infty[0,1]\), regarded as a
lattice-normed space over itself.  It is decomposable, \(bo\)-complete,
Dedekind complete, atomless, and has the strong Freudenthal property.

Let
\[
  F=\R,\qquad W=\Lone,\qquad q(g)=\|g\|_1.
\]
Every Banach space is a Banach--Kantorovich space over \(\R\) with its norm as
the vector norm.  Define
\[
  T:V\longrightarrow W,\qquad Tf=f.
\]
This is the natural inclusion \(L_\infty[0,1]\subset L_1[0,1]\).

\begin{theorem}
\label{thm:counterexample}
The operator \(T:\Linfty\to\Lone\) is a dominated linear operator whose exact
dominant is order narrow, but \(T\) is not order narrow.
\end{theorem}

\begin{proof}
For \(e\in E_+\), define
\[
  S(e)=\int_0^1 e(t)\,dt.
\]
Then for every \(f\in V\),
\[
  q(Tf)=\|f\|_1=\int_0^1 |f(t)|\,dt=S(p(f)),
\]
so \(S\) is a dominant of \(T\).

It is the exact dominant.  Indeed, if \(R:E\to\R\) is any positive dominant,
then for every \(e\in E_+\) we may take \(f=e\) and obtain
\[
  S(e)=\|e\|_1=q(Te)\leq R(p(e))=R(e).
\]
Thus \(S\leq R\), proving that \(S=|T|\).

Next \(S\) is order narrow.  Let \(e\in E\).  Write \(e=e^+-e^-\).  Since the
Lebesgue measure algebra is atomless, there are measurable sets
\(A_+\subset\{e>0\}\) and \(A_-\subset\{e<0\}\) such that
\[
  \int_{A_+} e^+\,dm=\frac12\int e^+\,dm,
  \qquad
  \int_{A_-} e^-\,dm=\frac12\int e^-\,dm.
\]
Put \(A=A_+\cup A_-\).  Then
\[
  e=e\chi_A\sqcup e\chi_{A^c}
\]
is a mutually complemented fragment decomposition in \(E\), and
\[
  S(e\chi_A)-S(e\chi_{A^c})=0.
\]
So the constant net given by this decomposition witnesses order narrowness of
\(S\) at \(e\).

Finally \(T\) is not order narrow.  Consider \(v=\one\in V\).  If
\[
  \one=v_1\sqcup v_2
\]
is a mutually complemented fragment decomposition in \(V\), then necessarily
\[
  v_1=\chi_A,\qquad v_2=\chi_{A^c}
\]
for some measurable \(A\subset[0,1]\).  Hence
\[
  q(Tv_1-Tv_2)
  =\|\chi_A-\chi_{A^c}\|_1
  =1.
\]
Thus no net of mutually complemented fragment decompositions of \(\one\) can
make \(q(Tv_1-Tv_2)\) converge to \(0\).  Therefore \(T\) is not order narrow.
\end{proof}

\section*{Why This Answers the Source Question}

The construction lies inside the structural setting of the theorem preceding
the source remark:
\[
  (V,E)=(L_\infty,L_\infty),\qquad (W,F)=(L_1,\R),
\]
with \(E\) atomless and Dedekind complete and \(F\) Dedekind complete.  The
exact dominant exists explicitly and is order narrow, yet the original
dominated linear operator fails order narrowness.  Hence the reverse
implication asked in the remark is false.

\section*{Novelty Notes}

The local run indexes were searched for the arXiv id \texttt{1508.03275}, the
paper title, and phrases around \texttt{exact dominant} and
\texttt{order narrow}.  No duplicate solution packet was found.  Local source
search found related vector-lattice and abstract-Uryson modulus results, but
those do not remove this lattice-normed-space counterexample.

Bounded web searches on July 1, 2026 for the exact question phrase, the paper
title, and close \texttt{order narrow exact dominant lattice-normed} phrases
found the source arXiv page and a related orthogonally additive paper, but no
exact published resolution of the reverse implication.

\begin{thebibliography}{9}

\bibitem{Source}
N. Abasov, A. E. M. Megahed, and M. Pliev,
\emph{Dominated operators from a lattice-normed space to a sequence vector
lattice},
arXiv:1508.03275.

\end{thebibliography}

\end{document}
